% Additive decomposition identities for Sparse Readout Prism.

\section{SRP Score Decomposition Identities}
\label{app:score-decomposition}
\label{app:interpreting-sparse-readout-prism-figures}
\label{app:additional-diagnostics}

This appendix states the decomposition identities used by the figures and control
tables and fixes the bookkeeping conventions for displaying a trained readout
SAE. The identities restate the local score decomposition introduced in
\S\ref{sec:method-local-score}. All reported scores are evaluated in the
model's raw readout space. If training used centering or row normalization, the
SAE coefficients, reconstructed rows, and residual terms are first mapped back
to that space.

\begin{table}[!tbp]
\caption{\textbf{Selected scores and their sparse-feature coefficients.}
Rows instantiate the selected readout direction \(q_\alpha\) and corresponding
direction coefficient \(\beta_i(\alpha)\). Barred \(w\) and \(z\) denote
uniform token-set averages, and for top competitor margins,
\(\bar{w}_{\pi,\mathcal{R}}=\sum_r\pi_r w_r\) and
\(\bar{z}_{\pi,\mathcal{R},i}=\sum_r\pi_r z_{r,i}\), where \(\pi_r\) are the
nonnegative competitor weights summing to one over the competitor set
\(\mathcal{R}\). Here \(A,B\in\mathcal{V}\) are single tokens,
\(\mathcal{A},\mathcal{B}\subset\mathcal{V}\) are token groups, and \(u\) is
the top-ranked token under the original LM head.}
\label{tab:method-special-cases}
\centering
\footnotesize
\setlength{\tabcolsep}{3pt}
\renewcommand{\arraystretch}{1.14}
\begin{tabularx}{\linewidth}{@{}>{\RaggedRight\arraybackslash}p{0.34\linewidth}Z Z@{}}
\toprule
Selected score family & \(q_\alpha\) & \(\beta_i(\alpha)\) \\
\midrule
\textbf{Token logit} & \(w_A\) & \(z_{A,i}\) \\
\addlinespace[1pt]
\textbf{Vocabulary-mean contrast} & \(w_A-\bar{w}_{\mathcal{V}}\) &
\(z_{A,i}-\bar{z}_{\mathcal{V},i}\) \\
\addlinespace[1pt]
\textbf{Logit difference} & \(w_A-w_B\) & \(z_{A,i}-z_{B,i}\) \\
\addlinespace[1pt]
\textbf{Group contrast} &
\(\bar{w}_{\mathcal{A}}-\bar{w}_{\mathcal{B}}\) &
\(\bar{z}_{\mathcal{A},i}-\bar{z}_{\mathcal{B},i}\) \\
\addlinespace[1pt]
\textbf{Top competitor margin} &
\(w_u-\bar{w}_{\pi,\mathcal{R}}\) &
\(z_{u,i}-\bar{z}_{\pi,\mathcal{R},i}\) \\
\bottomrule
\end{tabularx}
\end{table}

\subsection{Raw-Space Row Decomposition}
\label{app:score-decomposition-row-factorization}

Let \(w_v\in\R^{\dmodel}\) be the column-vector form of
unembedding row \(v\), following \S\ref{sec:method}. Under the row
preprocessing convention used by the reported analyses,
\[
    \begin{aligned}
    x_v&=\frac{w_v-\mu}{a_v}, \qquad a_v>0,\\
    x_v&=\sum_i z^{\mathrm{pre}}_{v,i}d^{\mathrm{pre}}_i+r^{\mathrm{pre}}_v,
    \end{aligned}
\]
with \(a_v=1\) for unnormalized rows. The raw-space decomposition then uses
\[
\begin{aligned}
    z_{v,i}&=a_v z^{\mathrm{pre}}_{v,i},
    & d_i&=d^{\mathrm{pre}}_i,\\
    r_v&=a_v r^{\mathrm{pre}}_v.
\end{aligned}
\]
It follows that
\[
    w_v=\mu+\sum_i z_{v,i}d_i+r_v .
\]
The displayed local score decomposition contains only this shared offset, SAE feature
terms, and the explicit residual term, with no hidden dense base term added.

\paragraph{Why the offset stays outside the dictionary.}
Nothing forbids folding the shared offset into the dictionary, since
\(\mu\) is a direction in \(\R^{\dmodel}\) like any other and the row
reconstruction could be written with one additional atom. We keep it
explicit because that atom would carry coefficient one on every row by
construction, so it would occupy one of the \(k\) active slots for every
token and spend sparse capacity on a quantity that is identical
everywhere. Its decoder direction would also be drawn toward the row mean,
away from structure that distinguishes rows. An explicit
centering term is the standard arrangement in sparse autoencoder practice
\citep{bricken2023monosemanticity,gao2024scaling}, and centering output
embeddings in particular is a documented stabilizer
\citep{stollenwerk2026output}. Holding \(\mu\) outside the sparse code
lets the features model deviations from the shared row mean, and it keeps
the offset a separately visible term in the identity for a selected score below,
where its cancellation becomes a property of the selected direction,
independent of which atoms happened to be learned.

\subsection{Selected-Score Identity}
\label{app:score-decomposition-selected-score}

For a hidden state \(h_\ell\), let
\(\rstate=T_\ell(h_\ell)\) be the corresponding state presented to
the LM head, following \S\ref{sec:method}. A selected readout direction is
a linear function of vocabulary rows,
\[
    q_\alpha = \sum_{v\in\mathcal{V}} \alpha_v w_v,
    \qquad
    s_\alpha(\rstate)=\rstate^\top q_\alpha .
\]
Define the direction coefficient, selected-direction residual, and SAE
feature contribution as
\[
\begin{aligned}
    \beta_i(\alpha)&=\sum_v \alpha_v z_{v,i},
    & r_\alpha&=\sum_v\alpha_v r_v,\\
    c_i(\rstate,\alpha)&=\beta_i(\alpha)\,p_i(\rstate),
\end{aligned}
\]
where \(p_i(\rstate)=\rstate^\top d_i\) is the projection
introduced in \S\ref{sec:method-factorization}.
Substituting the raw-space row decomposition gives the exact grouped identity
\[
\begin{aligned}
    s_\alpha(\rstate)
    &=
    \smu(\rstate,\alpha)
    +\sfeat(\rstate,\alpha)\\
    &\quad
    +\sresid(\rstate,\alpha),\\
    \smu(\rstate,\alpha)
    &=
    \left(\sum_v \alpha_v\right)\rstate^\top\mu,\\
    \sfeat(\rstate,\alpha)
    &=
    \sum_i c_i(\rstate,\alpha),\\
    \sresid(\rstate,\alpha)
    &=
    \rstate^\top r_\alpha .
\end{aligned}
\]
For zero-sum contrasts, including pairwise logit differences, group contrasts,
vocabulary mean contrasts, and top competitor margins, the offset term
cancels. \Cref{fig:method-schematic} depicts the decomposition.

\subsection{Error, Sign, and Display Rules}
\label{app:score-decomposition-fidelity}
\label{app:display-taxonomy}
\label{app:display-residual-decomposition}

For a selected readout direction \(q_\alpha\), the exact and reconstructed
scores are
\[
\begin{aligned}
    \sexact
    &=
    \rstate^\top q_\alpha,\\
    \srecon
    &=
    \smu(\rstate,\alpha)
    +
    \sfeat(\rstate,\alpha).
\end{aligned}
\]
The residual score is the reconstruction error,
\[
    \epsilon(\rstate,\alpha)
    =
    \sexact-\srecon
    =
    \sresid(\rstate,\alpha),
\]
and aggregate summaries use the parameterized relative error
\[
    \rho_\delta(\rstate,\alpha)
    =
    \frac{|\epsilon(\rstate,\alpha)|}{|\sexact|+\delta},
\]
with \(\delta=0.5\) logits, denoted \(\rho_{0.5}\).
We drop the arguments of \(\epsilon\) and \(\rho_\delta\) where the state and
coefficient vector are clear from context.
The floor prevents arbitrarily small margins from
dominating median error summaries, and near the decision boundary sign agreement and
margin-binned flip rates are the primary reliability checks. The unfloored
variant
\(\rho_0(\rstate,\alpha)=|\epsilon(\rstate,\alpha)|/|\sexact|\)
is used for identity checks, case-study displays, and all
baseline comparisons. It upper-bounds \(\rho_{0.5}\), and the two subscripts are
never mixed within one table. For signed
contrasts, side-specific feature support is reported together with the sign check
\[
    \operatorname{sign}(\sexact)
    =
    \operatorname{sign}(\srecon).
\]
Each interpreted figure states the selected readout score, exact score,
reconstructed score, reconstruction error, relative reconstruction error, and
sign status for contrasts. If a plot displays only the largest bars, the
reported feature sum and residual are still computed from the full active
feature sum. Token-level figures report tokenization or row audit
information where it affects interpretation, and brittle single-token contrasts are
handled as explicit group contrasts when needed.
